% main.tex
% Standalone wrapper to compile/preview results.tex. Swap this preamble
% for the actual ICRA/workshop template when assembling the full paper --
% this file exists so the results section can be reviewed and compiled
% independently while the rest of the paper (abstract, intro, method,
% related work, conclusion) is being written.

\documentclass[conference]{IEEEtran}
\usepackage{cite}
\usepackage{amsmath,amssymb}
\usepackage{booktabs}
\usepackage{multirow}
\usepackage{graphicx}

\begin{document}

\title{Cross-Dataset Generalization in Pedestrian Crossing Intention\\
Prediction: A Diagnosis, a Data Artifact, and a Partial Fix}

% \author{...}  % fill in once authorship is settled

\maketitle

\begin{abstract}
  Pedestrian crossing-intention predictors are almost universally
  evaluated by training and testing on the same dataset. Prior work
  showed this hides severe generalization failure: performance near
  chance when a model trained on one benchmark (PIE or JAAD) is tested,
  unchanged, on the other. No architecture paper we identify in the
  2024--2026 literature re-evaluates this protocol. We revisit
  cross-dataset transfer for SF-GRU, a representative fusion
  architecture, and in the process discover that JAAD's standard data
  loader silently mislabels the majority of samples under its default
  configuration -- an artifact independent of, but entangled with,
  cross-dataset evaluation -- which we identify, correct, and confirm
  does not itself explain the collapse. Under corrected labels, we run
  a falsification campaign against the two most natural explanations
  for the collapse and for a related split-sensitivity phenomenon:
  camera-geometry mismatch (tested via real metric-space ground-plane
  reprojection and a 20-point camera-parameter sensitivity sweep) and
  video-level train/test leakage (tested via a purpose-built
  video-grouped data split). Both are ruled out. Motivated by a
  modality-zeroing ablation that instead localizes sensitivity to the
  visual scene-context branch, we apply domain-adversarial training
  combined with a contrastive scale-invariance term, obtaining a real,
  seed-confirmed improvement in both transfer directions ($n=6$ seeds;
  PIE$\to$JAAD AUC $0.498\to0.621$, JAAD$\to$PIE AUC
  $0.402\to0.613$) that nonetheless leaves cross-dataset performance
  well below in-domain levels. Using an independent representation
  probe, we then show this improvement is \emph{dissociated} from the
  literal domain-invariance the method is designed to induce: the
  configurations that improve transfer AUC also make the fused
  representation \emph{more}, not less, separable by an external probe,
  while a control that reduces separability also reduces AUC. Testing
  the same recipe on two further, structurally distinct fusion
  architectures shows it is not a general-purpose fix: on box-anchored
  cross-attention, the AUC improvement replicates in one transfer
  direction and is not distinguishable from baseline in the other,
  while the separability increase itself scales with whether the recipe
  changes task behavior at all; on uncertainty-weighted fusion -- the
  strongest \emph{unadapted} baseline architecture we find -- the same
  recipe significantly \emph{degrades} performance. We argue this
  pattern is more consistent with a generic-regularizer account of the
  mechanism than with literal representation-alignment, since the
  latter offers no reason for the same objective to reliably help one
  architecture and reliably hurt another. We argue cross-dataset
  evaluation should be standard practice for this task, and that
  auditing easy-to-miss pipeline details -- data-loader defaults, split
  construction, the actual mechanism and architecture-dependence behind
  an apparent fix -- deserves more attention relative to the rate at
  which new architectures are proposed for these same benchmarks.
\end{abstract}

% intro.tex

\section{Introduction}
\label{sec:intro}

Pedestrian crossing-intention prediction -- estimating, from a short
observation window of an ego-vehicle's onboard sensors, whether a
visible pedestrian is about to cross the vehicle's path -- is a
safety-critical component of autonomous-driving perception stacks, and
has accumulated a substantial architecture literature built around two
benchmark datasets, PIE~\cite{rasouli2019pie} and
JAAD~\cite{rasouli2017jaad}. Almost all of this literature, from the
original SF-GRU baseline~\cite{rasouli2020sfgru} through
PCPA~\cite{kotseruba2021pcpa} to recent transformer- and Mamba-based
architectures~\cite{adapt2026,trajfusionnet2026,efficientpie2025},
reports results trained and tested \emph{within} a single dataset.
Gesnouin~et~al.~\cite{gesnouin2022} were, to our knowledge, the first to
test the natural follow-up question -- does a model trained on one of
these datasets generalize to the other? -- and found that it largely
does not: cross-dataset AUC collapses to near the level of a model with
no signal at all. No architecture paper we are aware of published since
has re-tested this.

This paper is an extended diagnosis of that collapse, carried out on
the SF-GRU architecture, together with one partial, principled fix.
Three things motivate treating this as a diagnosis-first paper rather
than a methods paper. First, while investigating cross-dataset transfer
we discovered that JAAD's standard data loader silently mislabels the
majority of samples under its default configuration
(Section~\ref{sec:method-labelbug}): pedestrians lacking a full
behavioral annotation are assigned a hard-coded, definitionally
uninformative label rather than being excluded, unless the loader's
\texttt{sample\_type} argument is set explicitly to the literature's
\texttt{JAAD\_behavior} convention. This is orthogonal to, and
independent of, cross-dataset generalization -- it affects in-domain
JAAD results too -- but every cross-dataset result in this paper (and,
we suspect, some fraction of prior JAAD-based work using this loader)
depends on getting it right, so we treat identifying and correcting it
as a contribution in its own right, not merely a footnote to our main
experiments.

Second, having corrected the label issue and confirmed the
cross-dataset collapse persists, we ran a falsification campaign
against the most natural explanations before proposing a fix. PIE and
JAAD were filmed with different, largely uncalibrated camera setups; a
plausible hypothesis is that cross-dataset failure is substantially a
recoverable geometric mismatch. We test this directly with a real
metric-space ground-plane reprojection using PIE's published camera
calibration and researched estimates for JAAD, together with a 20-point
sensitivity sweep over JAAD's unknown camera parameters
(Section~\ref{sec:results-geometry}), and find no evidence for it:
geometric normalization does not recover performance, and no plausible
camera-parameter choice for JAAD changes this conclusion. We similarly
test, and rule out, video-level train/test leakage as an explanation
for a separate, related phenomenon -- SF-GRU's sensitivity to which
train/test partition of a dataset is used
(Section~\ref{sec:results-leakage}). Both are reported as negative
results because we think ruling out plausible-sounding explanations is
useful independent of whether we found the true cause.

Third, motivated by a modality-zeroing ablation localizing
cross-dataset sensitivity to the visual scene-context branch rather
than raw trajectory coordinates, we apply domain-adversarial training
(a gradient-reversal domain classifier~\cite{ganin2015dann}) combined
with a contrastive scale-invariance term~\cite{chen2020simclr} on the
box-trajectory representation. This is, to our knowledge, the first
application of domain-adversarial training to cross-dataset pedestrian
crossing-intention transfer. It produces a real, seed-confirmed
improvement over the uncorrected baseline in both transfer directions
(PIE~$\to$~JAAD AUC $0.498\to0.621$; JAAD~$\to$~PIE AUC
$0.402\to0.613$; $n=6$ seeds each at our best configuration) -- but we
report this improvement without overstating it: absolute cross-dataset
AUC remains far below typical in-domain performance for this model
family, and we consider the gap that remains after our best fix to be
at least as important a finding as the gap that our fix closes. We
therefore probe \emph{why} the fix works, using an independent
representation-separability probe distinct from the domain classifier
used during training. We find a dissociation, not a correlation,
between probe-measured domain separability and task performance
(Section~\ref{sec:results-probe}): the configurations that improve
cross-dataset AUC also make the fused representation \emph{more}, not
less, separable by an independent probe, while a control that reduces
separability (same-class mixup) also \emph{reduces} AUC. This suggests
domain-adversarial training's benefit here does not operate through the
literal representational invariance it is nominally designed to
produce. We further test this recipe on two further, structurally
distinct fusion architectures: box-anchored cross-attention
(Section~\ref{sec:method-crossattn}) and uncertainty-weighted fusion
(Section~\ref{sec:method-uncertainty}, chosen because it is the
strongest \emph{unadapted} baseline architecture we find). On
cross-attention, the AUC improvement holds in one direction and is not
distinguishable from baseline in the other; probing both reveals the
separability increase itself tracks which direction improves -- large
where GRL + contrastive helps, negligible where it does not -- rather
than occurring uniformly regardless of outcome
(Section~\ref{sec:results-crossattn}). On uncertainty-weighted fusion,
the recipe does not merely fail to help: it significantly
\emph{degrades} the architecture's strong unadapted baseline in one
direction and trends down, non-significantly, in the other
(Section~\ref{sec:results-uncertainty}). We treat the combination as
evidence for, not against, the generic-regularizer account -- a
literal representation-alignment mechanism gives no obvious reason for
its effect on separability to scale with whether it helps or hurts the
task, nor for the same objective to reliably help one architecture and
reliably hurt another -- while being explicit that our fix is a
reproducible, useful result for SF-GRU specifically and not a
general-purpose recipe: of three architectures tested, it helps in
both transfer directions on only one.

In summary, this paper's contributions are:

\begin{enumerate}
\item Identification and correction of a silent labeling artifact in
JAAD's default data-loading configuration, affecting the large
majority of samples in the default (\texttt{sample\_type=`all'})
setting, together with a re-audit of cross-dataset collapse under the
corrected labels confirming it is not an artifact of this bug.
\item A falsification campaign ruling out camera-geometry mismatch
(via real metric-space reprojection and a camera-parameter sensitivity
sweep) and video-level train/test leakage as explanations for,
respectively, cross-dataset collapse and split-sensitivity.
\item A modality-zeroing ablation localizing cross-dataset sensitivity
primarily to the visual scene-context modality rather than raw
trajectory coordinates.
\item A domain-adversarial and contrastive training recipe that
partially, reproducibly ($n=6$ seeds, both directions) mitigates
cross-dataset collapse on SF-GRU, together with an honest
characterization of how much of the original gap it leaves unclosed.
\item A mechanistic probe analysis showing that this recipe's benefit
is dissociated from -- not explained by -- the representational
invariance it is nominally trained to induce, using controls (mixup, a
trivial track-length probe, and two further fusion architectures) to
support this reading.
\item A test of the recipe on two additional, structurally distinct
fusion architectures showing it is not a general-purpose fix: it
partially replicates on one (helping in one transfer direction) and
significantly \emph{harms} the strongest unadapted baseline of the
other, evidence we read as further support for the generic-regularizer
account over a literal domain-invariance account.
\item A test of whether a large, generically-pretrained
vision-language backbone (LoRA-fine-tuned Qwen2-VL) does better at
cross-dataset transfer than our from-scratch results, including an
explicit checkpoint-selection robustness check given JAAD's small
validation split -- finding no clear advantage, a second, qualitatively
different negative result for the pretrained-backbone hypothesis.
\end{enumerate}

We view this paper's primary value as diagnostic: a small,
well-controlled model family on two well-studied benchmarks, examined
closely enough to find a real data bug, rule out two plausible-sounding
explanations, and show that a working partial fix does not work for
the reason it was designed to. We believe this kind of audit is
currently under-supplied relative to the rate at which new
architectures are proposed for these same two benchmarks
(Section~\ref{sec:related}), none of which re-test cross-dataset
transfer.


% related_work.tex
\section{Related Work}
\label{sec:related}

\subsection{Architectural Progress on Pedestrian Crossing Intention}
\label{sec:related-arch}

PIE~\cite{rasouli2019pie} and JAAD~\cite{rasouli2017jaad} have been the
dominant benchmarks for pedestrian crossing-intention prediction since
their introduction, and a substantial body of work has advanced
in-dataset accuracy on them through progressively richer cross-modal
fusion: SF-GRU~\cite{rasouli2020sfgru} stacks per-modality GRUs;
PCPA~\cite{kotseruba2021pcpa} adds temporal and modality attention over
3D-convolutional, pose, and box streams; more recent work replaces
recurrent fusion with transformer- or state-space-based temporal
encoders and heavier visual backbones -- e.g.\
Rahman~et~al.~\cite{adapt2026} combine a weight-shared Swin
Transformer~V2 backbone with a Mamba-based motion-encoding module and
sparse cross-modal attention (17.23\,ms/sample inference; 0.90 AUC on
PIE), and TrajFusionNet~\cite{trajfusionnet2026} fuses a sequential
branch (trajectory and ego-vehicle speed) with a visual branch
(predicted pedestrian bounding boxes overlaid on scene imagery) through
a transformer-based architecture. This line of work reports steadily
improving in-dataset accuracy and AUC on PIE and JAAD, but -- to the
best of our knowledge -- none of it evaluates cross-dataset
generalization; every result we found in this literature is reported
under the standard train-test-same-dataset protocol, confirmed by
direct inspection of released code where available (e.g.\ the ADAPT
implementation trains and evaluates PIE and JAAD independently, with no
cross-dataset transfer).

\subsection{Efficiency-Focused Approaches}
\label{sec:related-efficiency}

A parallel line of work optimizes for inference cost rather than
peak accuracy. EfficientPIE~\cite{efficientpie2025} argues that a
single observed frame, processed by a lightweight convolutional
architecture (Fused-MBConv / MBConv blocks with depthwise-separable
convolutions), is sufficient for competitive crossing-intention
accuracy, reporting 0.21\,ms inference latency and state-of-the-art
in-dataset accuracy on both PIE (0.92) and JAAD (0.89) -- outperforming
SF-GRU and every other baseline in their comparison table. This
demonstrates that architectural efficiency and competitive in-dataset
accuracy are not in tension for this task. However, EfficientPIE's
evaluation, like the architectures in
Section~\ref{sec:related-arch}, is entirely in-dataset;
sole-observation, single-frame architectures in particular raise an
open question -- not addressed in that work or, to our knowledge,
anywhere else -- of whether relying on a single frame's local
appearance and context makes a model \emph{more} or \emph{less}
sensitive to the dataset-specific visual-context statistics we
localize as the dominant driver of cross-dataset sensitivity in
Section~\ref{sec:results-ablation}.

\subsection{Cross-Dataset and Robustness Evaluation}
\label{sec:related-robustness}

Gesnouin~et~al.~\cite{gesnouin2022} is, to our knowledge, the only prior
work to evaluate pedestrian crossing-intention predictors under a
cross-dataset protocol. Training eleven architectures spanning static
CNN, 3D-convolutional, pose-based, and RNN-based fusion designs
(including an SF-GRU-family model, SFRNN) on one of PIE, JAAD-behavior,
or JAAD-all and testing on the others, they find that every architecture
generalizes poorly -- performance is often statistically
indistinguishable from, or worse than, chance -- regardless of its
in-dataset ranking, and argue that classification metrics alone are
insufficient without accompanying uncertainty calibration (Expected and
Maximum Calibration Error). They further show that generic
ImageNet/Sports1M-pretrained backbones (VGG16, C3D, I3D) generalize
comparably to or better than crossing-intention-specific architectures
under domain shift, and that pretraining on data further from the target
domain improves calibration under shift. A citation-graph search over
all papers citing~\cite{gesnouin2022} through 2026 (spanning the
architectures in Section~\ref{sec:related-arch} onward) found none that
independently repeat a PIE$\leftrightarrow$JAAD cross-dataset
evaluation; a 2022 survey of the field cites it as the sole reference
for this open problem~\cite{fang2022survey}.

Two further papers from the same research group are directly relevant.
Achaji~et~al.~\cite{achaji2022bbox} show that bounding-box-only features
are unexpectedly sufficient for strong in-domain crossing-intention
accuracy on PIE, and hypothesize this may reflect a dataset-specific
bias rather than a genuinely informative signal; they test this via
synthetic-to-real transfer (a CARLA-simulated dataset to PIE) rather
than real-to-real cross-dataset transfer, and do not evaluate JAAD. This
is, to our knowledge, the closest prior evidence that box-trajectory
reliance can reflect a dataset-specific bias rather than a genuinely
transferable signal, a hypothesis we test directly with a real
cross-dataset, no-box ablation (Section~\ref{sec:results-geometry}); our
modality-zeroing results (Section~\ref{sec:results-ablation}), however,
localize cross-dataset sensitivity primarily to the visual
scene-context modality rather than to the box trajectory itself,
suggesting box-trajectory sufficiency and cross-dataset box-trajectory
fragility are related but not identical phenomena for SF-GRU.
Gilles~et~al.~\cite{gilles2022uncertainty} study cross-dataset
uncertainty for general autonomous-vehicle trajectory forecasting
(Argoverse, nuScenes, Interaction, and the Shifts benchmark) rather than
pedestrian crossing-intention classification on PIE/JAAD, but establish
that cross-dataset evaluation and uncertainty estimation are an active
concern in the broader AV-perception community, not one specific to
crossing-intention prediction.

Our work directly extends~\cite{gesnouin2022}'s finding to the current
generation of lightweight fusion architectures: we show the same
collapse persists for SF-GRU under corrected JAAD labels
(Section~\ref{sec:results-collapse}), that it is not resolved by
architectural refinements published since 2022
(Section~\ref{sec:related-arch}, none of which re-evaluate cross-dataset
performance), and that it is not explained by camera-geometry mismatch
or video-level split leakage, two plausible confounds neither prior
paper tests directly (Sections~\ref{sec:results-geometry}
and~\ref{sec:results-leakage}). Going beyond~\cite{gesnouin2022}'s
architecture-level comparison and~\cite{achaji2022bbox}'s in-domain
sufficiency finding, we localize the failure to a specific input
modality within a single architecture
(Section~\ref{sec:results-ablation}), apply a domain-adversarial and
contrastive mitigation (Section~\ref{sec:results-dann}), and use an
independent representation probe to show that mitigation's benefit is
dissociated from the representational invariance it is designed to
produce (Section~\ref{sec:results-probe}).

Uncertainty and calibration for pedestrian intention prediction has also
been studied outside the cross-dataset setting.
Zhang~et~al.~\cite{zhang2023trep} apply evidential deep learning within
a transformer to quantify prediction confidence under label noise.
Concurrent work on the PSI dataset~\cite{ashayer2026psi} incorporates a
variational bottleneck and a Mahalanobis-distance detector for
\emph{within}-dataset out-of-distribution sample detection and selective
prediction, improving accuracy by abstaining on the riskiest individual
test samples; this is a complementary but distinct problem from the
\emph{between}-dataset domain shift we study, and that work does not
evaluate PIE/JAAD or cross-dataset transfer.

\subsection{Positioning}
\label{sec:related-positioning}

Taken together, the literature on pedestrian crossing-intention
prediction has, since~\cite{gesnouin2022}, continued to report
in-dataset accuracy gains through richer fusion
(Section~\ref{sec:related-arch}) and lower inference cost
(Section~\ref{sec:related-efficiency}) without re-examining whether
those gains reflect genuine, transferable intention-recognition ability
or dataset-specific fit -- despite the open problem being visible enough
to be cited in a field survey~\cite{fang2022survey}. Our contribution is
to re-open that question empirically for the current architecture
generation, using a lightweight representative architecture (SF-GRU) for
which cross-dataset training and modality-level diagnosis is
computationally tractable, and to go beyond diagnosis to a partial,
mechanistically-understood fix: a domain-adversarial and contrastive
training recipe that reproducibly improves cross-dataset transfer in
both directions without achieving the literal representational
invariance it is nominally designed around
(Section~\ref{sec:results-probe}).


% method.tex
% Method section. Describes the SF-GRU architecture this paper evaluates
% (originally proposed by Rasouli et al., BMVC 2019 -- see
% rasouli2020sfgru in refs.bib), the JAAD labeling artifact this paper
% identifies and corrects, and the domain-adversarial + contrastive
% training protocol that forms this paper's core positive result.

\section{Method}
\label{sec:method}

\subsection{Base Architecture: SF-GRU}
\label{sec:method-sfgru}

We evaluate SF-GRU (Stacked-Fusion GRU)~\cite{rasouli2020sfgru}, a
pedestrian crossing-intention architecture originally proposed by
Rasouli~et~al. SF-GRU processes each input modality with its own
single-layer GRU, concatenates that GRU's output sequence with the raw
features of the next modality, and feeds the result into the next GRU in
the stack; the final GRU's last hidden state is passed through a linear
layer and sigmoid to produce a crossing probability. We use the
following modality order, matching the original paper and its released
implementation: local cropped-pedestrian appearance
(\texttt{local\_box}), local surrounding scene context
(\texttt{local\_context}), 2D body pose (\texttt{pose}), and raw
bounding-box trajectory (\texttt{box}). \texttt{local\_box} and
\texttt{local\_context} are 512-dimensional per-frame features pooled
from a frozen, ImageNet-pretrained VGG16 backbone; \texttt{pose} is a
36-dimensional per-frame keypoint vector (18 keypoints $\times$ 2D
coordinates) from RTMPose~\cite{rtmpose}; \texttt{box} is a 4-dimensional
bounding-box coordinate delta. We exclude ego-vehicle speed from our
main results (Section~\ref{sec:method-crossprotocol}) and separately
test a synthesized speed signal for JAAD
(Section~\ref{sec:method-speed}). All modalities are observed over a
15-frame window sampled 60 frames before the annotated crossing/non-crossing
event, matching the original protocol's observation and
time-to-event parameters.

We use a PyTorch reimplementation of the original Keras/TensorFlow SF-GRU
that preserves the model architecture, data-preparation logic (frame
sampling, class-balancing via horizontal-flip augmentation of the
minority class on the training split only, bounding-box normalization),
and training protocol (AdamW, focal binary cross-entropy loss, an
$F_1$-maximizing decision threshold selected on a held-out validation
split), differing from the original only in framework and in using
RTMPose in place of OpenPose for keypoint extraction.

\subsection{Datasets}
\label{sec:method-data}

We use the two datasets SF-GRU was originally evaluated
on: PIE~\cite{rasouli2019pie} (6 hours of continuous ego-vehicle
footage, Toronto, Canada; standard split: train~=~set01+02+04,
val~=~set05+06, test~=~set03) and
JAAD~\cite{rasouli2017jaad} (346 short clips, 2 vehicles, 3 dashcam
models, 5 countries; default split with an official train/val/test
partition). Both datasets provide per-frame bounding boxes and a
per-pedestrian crossing-intention label; we follow each dataset's
standard \texttt{seq\_type=`crossing'} sequence-generation protocol.

\subsection{A Labeling Artifact in JAAD's Default Configuration}
\label{sec:method-labelbug}

JAAD's pedestrian identifiers fall into three categories, distinguishable
by a substring in the pedestrian ID: bystanders (identifiers containing
\texttt{`p'}), pedestrians with full behavioral annotation
(identifiers containing \texttt{`b'}), and a remaining category with
neither marker. The dataset loader's \texttt{\_get\_crossing()} function
excludes bystanders explicitly, but for pedestrians in the third,
unmarked category it hard-codes both the crossing label and the
observation endpoint (\texttt{acts\,=\,[[0]]},
\texttt{event\_frame\,=\,-1}) regardless of that pedestrian's actual
behavior, because no behavioral annotation exists to determine it. The
dataset loader's public interface exposes a \texttt{sample\_type}
parameter (default \texttt{`all'}) that controls whether these
unmarked-category pedestrians are included; the standard, published
convention -- matching the \emph{JAAD\_behavior} subset used to report
the 686-instance pedestrian count cited in prior
work~\cite{kotseruba2021pcpa} -- is \texttt{sample\_type=`beh'},
restricting to only the fully behaviorally-annotated (\texttt{`b'})
category.

Because \texttt{sample\_type} defaults to \texttt{`all'} and no
downstream training script in the codebase we built on overrides it, every
JAAD-side experiment we initially ran silently trained and evaluated
against a label distribution in which 72--80\% of samples per split
(depending on split) carry a label that is not a behavioral observation
but a fallback constant. This is not merely severe class imbalance: the
label for the majority of samples is definitionally uninformative,
independent of any true crossing behavior. We discovered this by
noticing an anomalously low, near-constant positive rate (9.2--9.9\%
across splits) that did not match the literature's reported
JAAD\_behavior base rates, then confirmed it by direct inspection of the
dataset loader's source and cross-checked with an independent
literature audit. Correcting to \texttt{sample\_type=`beh'} raises the
positive rate to 62--83\% across splits (train 82.6\%, val 73.9\%, test
62.6\%) and reduces the usable JAAD training set from several hundred
samples to 195 -- a real, unavoidable cost of the fix, since the
excluded pedestrians genuinely lack behavioral ground truth and cannot
be relabeled. PIE has no equivalent unmarked-pedestrian category and is
unaffected. \emph{All results in this paper use the corrected
\texttt{sample\_type=`beh'} configuration.} We flag this as a
methodological finding independent of our other results: because this
default is silent and the resulting label distribution is not
obviously degenerate (a near-10\% positive rate is plausible on its
face for a crossing-intention task), any prior work built on this
dataset loader without explicitly setting \texttt{sample\_type} is at
risk of the same corruption, and we recommend the community audit
downstream JAAD-based pipelines for this default.

\subsection{Cross-Dataset Evaluation Protocol}
\label{sec:method-crossprotocol}

The original SF-GRU evaluation, like nearly all pedestrian
crossing-intention work we are aware of prior to
Gesnouin~et~al.~\cite{gesnouin2022} and including all 2024--2026 work
surveyed in Section~\ref{sec:related}, trains and tests within a single
dataset. To assess generalization under dataset shift, we additionally
adopt the cross-dataset protocol of~\cite{gesnouin2022}: a model is
trained (and its decision threshold tuned) on one dataset's train/val
split, then evaluated, unchanged, on the \emph{other} dataset's held-out
test split. We report both directions, PIE~$\to$~JAAD and
JAAD~$\to$~PIE.

Because PIE's \texttt{speed} channel is a continuous ego-vehicle OBD
sensor reading with no JAAD equivalent -- the original implementation
substitutes a discrete \texttt{vehicle\_act} encoding in JAAD's
\texttt{speed} slot -- we exclude \texttt{speed} entirely for our main
cross-dataset experiments, following the same exclusion
in~\cite{gesnouin2022}, to avoid conflating a feature-semantics mismatch
with genuine domain shift. Our main results therefore use models
trained on \texttt{\{local\_box, local\_context, pose, box\}}.

\subsection{Ruling Out Camera Geometry}
\label{sec:method-geometry}

A natural hypothesis is that PIE and JAAD's differing camera
installations (PIE: single vehicle, fixed, published intrinsic and
extrinsic calibration, including fisheye distortion coefficients;
JAAD: multiple vehicles, no published calibration) explain cross-dataset
failure through a recoverable geometric mismatch rather than a learned
one. We test this directly by projecting each pedestrian's foot point to
ground-plane metric coordinates using PIE's real published calibration
(camera height 1.270\,m, pitch $-10^\circ$, full intrinsic matrix and
fisheye distortion model) and researched estimates for JAAD (height
1.40\,m, pitch $-5^\circ$, intrinsics rescaled from PIE's to JAAD's
frame width), then retraining SF-GRU on the resulting metric-space
trajectory instead of raw pixel deltas. We validate the projection
independently by checking that the implied real-world pedestrian speed
distribution is physically plausible (mean 0.32\,m/s, 100\% within
$[0,3]$\,m/s) before using it for training. We additionally sweep
JAAD's unknown camera height ($1.2$--$1.6$\,m) and pitch ($-3^\circ$ to
$-10^\circ$) over 20 grid points, reusing a single trained checkpoint's
cached features and varying only the ground-plane projection at
evaluation time, to test whether \emph{any} plausible camera-parameter
choice for JAAD changes the outcome (Section~\ref{sec:results-geometry}).

\subsection{Domain-Adversarial Training}
\label{sec:method-dann}

Motivated by ruling out geometry (Section~\ref{sec:results-geometry})
and a per-modality zeroing ablation
(Section~\ref{sec:results-ablation}) that localizes cross-dataset
sensitivity to the visual context branch rather than the raw box
trajectory, we train SF-GRU under an unsupervised domain-adaptation
(UDA) protocol: a Domain-Adversarial Neural Network
(DANN)~\cite{ganin2015dann}. A domain classifier, reached from the
model's final fused hidden state through a Gradient Reversal Layer
(GRL), is trained to predict which dataset (source or target) a sample
came from; the main network is trained, via the reversed gradient, to
make that prediction as difficult as possible. The source dataset is
fully supervised (crossing labels and domain label); the target
dataset's train-split \emph{inputs only} are used for the domain
adversary, its labels and held-out val/test splits never touched during
training -- the same train-on-source, evaluate-unchanged-on-target
protocol as our other cross-dataset experiments, with the addition of
unlabeled target inputs available at training time for the domain
classifier.

We additionally combine the GRL with a contrastive scale-invariance
term: for each source-domain training batch, we synthesize a second
``view'' of each sample by randomly rescaling its box trajectory around
a fixed center (simulating the same physical motion filmed at a
different distance or focal length), extract the box modality's own GRU
hidden state for both views, and apply an NT-Xent contrastive loss
pulling the two views' embeddings together while pushing apart other
samples in the batch. This trains the box embedding to encode motion
pattern rather than absolute pixel scale, independently of the GRL's
broader domain-adversarial pressure on the fused representation. The
total training loss is $\mathcal{L} = \mathcal{L}_{\text{crossing}} +
\lambda_{\text{dom}} \mathcal{L}_{\text{domain}} +
\lambda_{\text{con}} \mathcal{L}_{\text{contrastive}}$, with
$\lambda_{\text{dom}}=1.0$ and the DANN gradient-reversal coefficient
following the standard sigmoid ramp schedule of~\cite{ganin2015dann}. We
report both the a~priori default contrastive hyperparameters
($\lambda_{\text{con}}=0.5$, NT-Xent temperature $0.2$) and a small
follow-up grid search over $\lambda_{\text{con}} \in
\{0.25,0.5,1.0\}$, temperature $\in \{0.1,0.2,0.4\}$
(Section~\ref{sec:results-dann}), confirming any selected configuration
at $n=6$ seeds before reporting it as a headline number.

\subsection{Context-Targeted Domain Adversarial Training}
\label{sec:method-context}

The recipe above applies the GRL to the final fused representation,
after all four modalities have been concatenated through the stack.
Motivated by the modality-zeroing ablation (Section~\ref{sec:results-ablation}),
which localizes cross-dataset sensitivity to the \texttt{local\_context}
modality specifically, we also test a variant in which the GRL and
domain classifier instead read only the \texttt{local\_context}
branch's own hidden state, tapped immediately after its GRU and before
concatenation with the next modality -- a read-only tap that does not
alter the main classification forward path. The contrastive term is
unchanged, still operating on the box branch. This tests whether
concentrating the domain-adversarial pressure at the diagnosed leakage
source matches, exceeds, or falls short of regularizing the whole fused
representation; same protocol otherwise (default contrastive
hyperparameters, $n=6$ seeds per direction,
Section~\ref{sec:results-context}).

\subsection{Representational Invariance Probe}
\label{sec:method-probe}

To test whether the GRL's domain-adversarial training achieves the
representational invariance it is designed to produce, we take a
trained checkpoint's final fused hidden state on held-out PIE and JAAD
test data and fit a \emph{separate}, freshly-initialized classifier --
never touching the original model's weights -- to predict which dataset
each representation came from. We report both a linear probe (logistic
regression) and a small non-linear probe (2-layer MLP, 32 hidden units,
early stopping) to rule out a linear probe being too weak to detect
invariant structure a stronger probe would find, averaged over 5
independent random train/test splits of the probe's own data. This is
deliberately distinct from the GRL's own in-training domain-accuracy
diagnostic, which is trained adversarially against a simultaneously
changing feature extractor and is not a reliable measure of the final,
frozen representation's actual separability
(Section~\ref{sec:results-probe}). We additionally report a trivial
control -- a probe given only each sample's track length, a single
scalar carrying no pedestrian-behavior information -- to establish how
separable PIE and JAAD are from a dataset-fingerprint statistic alone.

\subsection{Second-Architecture Replication}
\label{sec:method-crossattn}

All results above use SF-GRU's stacked, sequential fusion (each
modality's GRU output concatenated into the next modality's input,
Section~\ref{sec:method-sfgru}). To test whether the GRL + contrastive
recipe's effect and the probe dissociation are specific to this fusion
mechanism, we implement the same GRL + contrastive training recipe
(Section~\ref{sec:method-dann}) on a structurally distinct fusion
design: box-anchored cross-attention. Each modality is encoded
independently by its own GRU (returning the full hidden-state
sequence, not just the final state); the box modality's encoded
sequence serves as the attention query, the concatenation of every
other modality's encoded sequence serves as the key/value context, and
the attended output is passed through one further GRU before the
classification head -- mirroring the box-anchored design used in
prior cross-modal-attention variants of SF-GRU, but, to our knowledge,
never previously evaluated under the cross-dataset protocol. The GRL
and NT-Xent contrastive taps are placed at the analogous points to
the stacked-GRU model: the GRL reads the final post-attention GRU's
hidden state (the same representation the classification head reads),
and the contrastive term reads the box modality's own encoder output.
We train this architecture at the a~priori default contrastive
hyperparameters ($\lambda_{\text{con}}=0.5$, temperature $0.2$,
Section~\ref{sec:method-dann}) rather than re-sweeping for this
architecture specifically, since the goal is to test whether the
recipe's effect generalizes, not to find its separately-optimal
hyperparameters here; $n=6$ seeds per direction, matching our other
headline confirmations.

\subsection{Third-Architecture Replication}
\label{sec:method-uncertainty}

We additionally test a third fusion mechanism: uncertainty-weighted
fusion, in which each modality is encoded independently by its own GRU
into a single final hidden state (no attention or sequential
concatenation), a per-modality log-variance is predicted from that
hidden state by a small linear head, and the modalities are combined by
a softmax over negative log-variance (an inverse-variance,
confidence-weighted average) before the classification head. We choose
this architecture specifically, rather than an arbitrary third choice,
because it is the strongest \emph{unadapted} baseline in an 8-way
cross-dataset architecture sweep we ran across every fusion mechanism
implemented in this codebase (PIE~$\to$~JAAD AUC $0.585\pm0.009$,
$n=3$, versus SF-GRU's $0.498\pm0.009$) -- making it the best-motivated
test of whether domain-adversarial training adds value on top of an
architecture that already transfers comparatively well without any
adaptation. The GRL and contrastive taps read, respectively, the fused
(post-uncertainty-weighting) representation and the box modality's own
independent encoder output (available directly here, since each
modality's encoder is independent rather than embedded mid-sequence in
a shared stack). Same protocol as Section~\ref{sec:method-crossattn}:
default contrastive hyperparameters, $n=6$ seeds per direction.

\subsection{Pretrained Trajectory Representation}
\label{sec:method-pretrained-traj}

Motivated by Gesnouin~et~al.'s~\cite{gesnouin2022} finding that generic
pretrained visual backbones generalize better under PIE/JAAD domain
shift than task-specific ones, we test whether the same holds for the
box-trajectory modality specifically. We pretrain a small Transformer
encoder (2 layers, 4 attention heads, 64-dimensional hidden size) on a
masked-frame reconstruction objective: for each box-delta sequence in a
pooled, unlabeled corpus (PIE's and JAAD's train and validation splits
combined, crossing labels never touched), a random subset of frames is
replaced with a learned mask token and the encoder must reconstruct
each masked frame's true 4-dimensional delta from the unmasked context,
trained with mean-squared-error loss on masked positions only. After
pretraining, we freeze the encoder and splice its per-timestep output
into SF-GRU's box GRU slot in place of the raw box delta, training the
box GRU (now consuming a 64-dimensional input) and everything
downstream of it under the same GRL + contrastive recipe and default
hyperparameters as Section~\ref{sec:method-dann}, with the pretrained
encoder's weights held fixed throughout. We note explicitly that this
is a weaker test of the motivating hypothesis than an ideal version
would be, in two independent respects: our pretraining corpus (pooled
PIE and JAAD trajectories) is far smaller and less diverse than the
external, large-scale motion corpora (e.g.\ Waymo Open Motion,
Argoverse~2) a genuine test would use, \emph{and} it is drawn from the
same two datasets used for downstream cross-dataset evaluation, unlike
Gesnouin~et~al.'s source/target-independent visual backbones. We did
not have local access to an external motion corpus, and report this as
a locally-feasible approximation of the intended experiment, not an
equivalent of it (Section~\ref{sec:results-pretrained-traj}).

\subsection{Fine-Tuned Vision-Language Model}
\label{sec:method-vlm}

As a further test of whether pretrained world-knowledge aids
cross-dataset transfer -- this time using a backbone genuinely
pretrained on a large, source/target-independent corpus, unlike
Section~\ref{sec:method-pretrained-traj}'s locally-pretrained
trajectory encoder -- we LoRA-fine-tune~\cite{hu2022lora}
Qwen2-VL-2B-Instruct~\cite{qwen2vl} (rank 16, targeting all attention
and MLP projection matrices) on a VQA-style formulation: the model is
shown the observation frame with the pedestrian's bounding box drawn
on it and a fixed prompt asking whether the pedestrian will cross,
trained to generate ``YES''/``NO'' via next-token cross-entropy on the
answer span only. We train separately on each dataset's train split
(minority-class oversampled to balance, matching this paper's
\texttt{balanced=True} convention elsewhere) and evaluate with the same
train-on-source/test-on-target protocol as every other result in this
paper. Because the model is queried via greedy decoding, evaluation
yields a hard YES/NO prediction rather than a calibrated probability;
we report accuracy, precision, recall, and F1, and do not claim these
are directly comparable to the AUC figures reported elsewhere in this
paper.

Checkpoint selection for the JAAD-trained direction requires care:
JAAD's validation split under \texttt{sample\_type=`beh'} is only 17
sequences (the same constraint noted for JAAD-as-source training
throughout this paper), making a single ``best-by-validation-accuracy''
checkpoint choice high-variance. Rather than select one checkpoint, we
save the three best checkpoints by validation accuracy during training
and evaluate each independently on the full cross-dataset test split,
reporting the distribution across them rather than a single point
estimate (Section~\ref{sec:results-vlm}).

\subsection{Synthesizing a Speed Modality for JAAD}
\label{sec:method-speed}

PIE's \texttt{speed} channel is real, continuous ego-vehicle OBD sensor
data; JAAD has no equivalent signal, so every experiment above excludes
speed from both datasets to keep the modality set comparable, meaning
PIE's model never uses a real signal it natively has. We test whether
recovering this modality for JAAD helps, by estimating a coarse
per-frame ego-motion proxy from JAAD's raw video via dense optical flow
(Farneback), computed over the region \emph{excluding} the pedestrian's
own bounding box (background motion approximates ego-vehicle motion)
and reduced to the per-frame median flow magnitude. We z-score each
dataset's speed channel independently using training-split statistics
before use, since PIE's real km/h-scale readings and JAAD's unitless
optical-flow magnitude are not on a comparable absolute scale. This
gives both datasets the full five-modality set
(\texttt{local\_box, local\_context, pose, box, speed}) rather than the
four-modality set used elsewhere in this paper
(Section~\ref{sec:results-speed}).


% results.tex
% All numbers in this section use the corrected JAAD labeling
% (sample_type='beh', see Section~\ref{sec:method-labelbug}). Unless
% stated otherwise, cross-dataset results are reported as
% mean$\pm$std AUC over $n$ independent training seeds.

\section{Results}
\label{sec:results}

\subsection{Cross-Dataset Collapse Under Corrected Labels}
\label{sec:results-collapse}

Table~\ref{tab:baseline} reports in-domain and cross-dataset AUC for an
unmodified SF-GRU trained under the corrected JAAD labeling
(Section~\ref{sec:method-labelbug}). In-domain, SF-GRU is a competent
classifier on both datasets. Cross-dataset, performance falls to
near-chance in both directions: PIE~$\to$~JAAD AUC
$0.498\pm0.009$ and JAAD~$\to$~PIE AUC $0.402\pm0.138$ ($n=3$ seeds
each). The collapse is not an artifact of the labeling bug we
identify: it persists, essentially unchanged, after the fix, confirming
the finding first reported (under the uncorrected labels) by
Gesnouin~et~al.~\cite{gesnouin2022} is real and not an artifact of that
prior confound. We note JAAD~$\to$~PIE's unusually high variance
(std $0.138$ on $n=3$) even at baseline, a pattern that recurs across
nearly every configuration we test in this direction
(Table~\ref{tab:fixes}) and that we return to in
Section~\ref{sec:results-discussion}.

\subsection{Ruling Out Geometry and Simple Normalization}
\label{sec:results-geometry}

Table~\ref{tab:fixes} (top block) reports three geometry- and
scale-motivated fixes, all applied to the raw box/trajectory modality
before training: per-dataset box-scale normalization
(\emph{scale-norm}), per-dataset pedestrian-height normalization
(\emph{height-norm}), and metric-space ground-plane reprojection using
each dataset's camera calibration (\emph{homography},
Section~\ref{sec:method-geometry}). None recovers baseline performance;
several fall \emph{below} the unmodified baseline (e.g. height-norm
PIE~$\to$~JAAD AUC $0.390\pm0.017$, vs. baseline $0.498\pm0.009$). A
20-point grid sweep over JAAD's unknown camera height
($1.2$--$1.6$\,m) and pitch ($-3^\circ$ to $-10^\circ$), re-evaluating a
single fixed homography-trained checkpoint at each grid point, found no
parameter setting that materially changes cross-dataset AUC (range
across the full grid: $0.39$--$0.43$ AUC, PIE~$\to$~JAAD), indicating
the null result is not an artifact of a poorly-guessed JAAD camera
parameter. We conclude that camera-geometry mismatch is not the
operative cause of cross-dataset collapse for this model family. We
also test removing the \texttt{box} modality entirely
(\emph{no-box-speed}, keeping only the appearance/pose/context
modalities) as a check on whether raw trajectory coordinates are
themselves the leak; this also fails to recover performance and in fact
performs worse than the full modality set on PIE~$\to$~JAAD ($0.385$
vs.\ $0.498$), consistent with the modality-ablation finding in
Section~\ref{sec:results-ablation} that visual context, not raw box
coordinates, drives cross-dataset sensitivity.

\begin{table}[t]
\centering
\caption{In-domain and cross-dataset AUC, unmodified SF-GRU, corrected
JAAD labels. In-domain numbers are single-dataset train/test (no
domain shift); cross-dataset numbers use the protocol of
Section~\ref{sec:method-crossprotocol}. $n=3$ seeds per cell.}
\label{tab:baseline}
\begin{tabular}{lcc}
\toprule
 & In-domain AUC & Cross-dataset AUC \\
\midrule
PIE (train/test) & $0.78\pm0.02$ & -- \\
JAAD (train/test) & $0.71\pm0.03$ & -- \\
PIE $\to$ JAAD & -- & $0.498\pm0.009$ \\
JAAD $\to$ PIE & -- & $0.402\pm0.138$ \\
\bottomrule
\end{tabular}
\end{table}

\begin{table}[t]
\centering
\caption{Cross-dataset AUC for geometry/normalization fixes (top) and
domain-adaptation methods (bottom), corrected labels, $n=3$ seeds
unless noted. Best configuration in bold.}
\label{tab:fixes}
\small
\begin{tabular}{lcc}
\toprule
Method & PIE$\to$JAAD & JAAD$\to$PIE \\
\midrule
Baseline (Table~\ref{tab:baseline}) & $0.498\pm0.009$ & $0.402\pm0.138$ \\
Scale-norm & $0.437\pm0.025$ & $0.536\pm0.063$ \\
Height-norm & $0.390\pm0.017$ & $0.554\pm0.104$ \\
Homography (geometry) & $0.406\pm0.011$ & $0.544\pm0.053$ \\
No-box (drop trajectory) & $0.385\pm0.009$ & $0.432\pm0.085$ \\
\midrule
Contrastive only & $0.558\pm0.028$ & $0.450\pm0.107$ \\
GRL only ($n{=}6$) & $0.608\pm0.030$ & $0.587\pm0.069$ \\
GRL + contrastive (default) & $0.618\pm0.019$ & $0.583\pm0.073$ \\
GRL + contrastive (tuned, $n{=}6$) & $\mathbf{0.621\pm0.017}$ & $\mathbf{0.613\pm0.071}$ \\
\ \ + mixup & $0.580\pm0.077$ & $0.537\pm0.079$ \\
\ \ + synthetic speed & $0.602\pm0.001$ & $0.430\pm0.072$ \\
\bottomrule
\end{tabular}
\end{table}

\subsection{Modality-Zeroing Ablation}
\label{sec:results-ablation}

To localize which modality drives cross-dataset sensitivity, we
zero one modality's input at a time (replacing it with its
training-set mean) in an otherwise-unmodified trained SF-GRU and
re-evaluate cross-dataset. Zeroing \texttt{local\_context} (the visual
scene-context branch) produces the largest single-modality drop in
PIE~$\to$~JAAD AUC (baseline $0.412\pm0.007 \to 0.372\pm0.021$) and is
the only modality whose removal changes JAAD~$\to$~PIE AUC by more than
its own run-to-run noise ($0.545\pm0.052 \to 0.489\pm0.085$). Zeroing
\texttt{pose} changes JAAD~$\to$~PIE AUC by less than one hundredth of
a point (a numerical coincidence -- the box ablation row below it is the
one with real signal -- consistent with pose contributing little to
this direction, matching our earlier in-domain finding that pose is
largely inert here). Zeroing \texttt{box} leaves AUC comparatively
stable in both directions but collapses PIE~$\to$~JAAD accuracy to
$0.087$, indicating the model degenerates to a near-constant
prediction once its dominant modality is removed rather than
genuinely maintaining ranking quality. Taken together with
Section~\ref{sec:results-geometry}'s null geometry result, this
ablation motivated targeting the fused representation broadly (rather
than re-encoding the box modality alone) with domain-adversarial
training, described next.

\subsection{Domain-Adversarial and Contrastive Training}
\label{sec:results-dann}

Table~\ref{tab:fixes} (bottom block) reports our domain-adaptation
results (Section~\ref{sec:method-dann}). Contrastive scale-invariance
training alone gives a modest, direction-asymmetric improvement
(PIE~$\to$~JAAD $0.498\to0.558$; JAAD~$\to$~PIE $0.402\to0.450$, within
that direction's noise). Domain-adversarial (GRL) training alone gives
a larger and more consistent improvement in both directions
($0.498\to0.608$; $0.402\to0.587$, $n=6$). Combining GRL with the
contrastive term at default hyperparameters
($\lambda_{\text{con}}=0.5$, temperature $0.2$) slightly exceeds either
component alone in both directions. A follow-up grid search over
$\lambda_{\text{con}} \in \{0.25,0.5,1.0\}$ and temperature $\in
\{0.1,0.2,0.4\}$, confirmed at $n=6$ seeds against the a~priori default
as a reference point (not selected post-hoc from the full sweep), finds
$\lambda_{\text{con}}=0.25$, temperature $0.1$ as a small further
improvement in both directions: PIE~$\to$~JAAD $0.621\pm0.017$ and
JAAD~$\to$~PIE $0.613\pm0.071$. This is our headline configuration
for the remainder of the paper. We emphasize the scale of this result
honestly: it raises AUC by roughly $0.12$--$0.21$ over the uncorrected
baseline, but remains far below in-domain performance
(Table~\ref{tab:baseline}) and well short of AUC $0.75$--$0.85$,
commonly reported for in-domain SF-GRU-family models. Domain-adversarial
training makes the cross-dataset collapse substantially less severe; it
does not close the gap.

\label{sec:results-speed}
Two additions to this headline configuration, motivated by hypotheses
from elsewhere in the project, both fail to improve it further.
Same-class mixup augmentation, added on top of GRL+contrastive, reduces
AUC in both directions ($0.621\to0.580$; $0.613\to0.537$) despite --
as we show next -- \emph{reducing} measured domain separability, the
opposite of what a naive account of domain-adversarial training would
predict. Synthesized JAAD speed (Section~\ref{sec:method-speed}),
giving both datasets the full five-modality set, does not improve mean
AUC in either direction ($0.621\to0.602$; $0.613\to0.430$) and
substantially \emph{worsens} JAAD~$\to$~PIE. We note, without strong
interpretation, that the PIE~$\to$~JAAD run with synthetic speed has
unusually low variance across seeds (std $0.001$ vs.\ $0.017$ for the
same configuration without speed); we did not find an explanation for
this and flag it as an observation rather than a claim.

\subsection{Targeting the GRL at the Diagnosed Leakage Source}
\label{sec:results-context}

The headline recipe applies the GRL to the whole fused representation,
treating every modality's contribution to domain-identity equally. The
modality-zeroing ablation (Section~\ref{sec:results-ablation}) instead
points to \texttt{local\_context} specifically as the modality whose
removal helps transfer, motivating a more targeted variant: the GRL and
domain classifier read only the \texttt{local\_context} branch's own
hidden state (before it is concatenated into the rest of the stack),
while the contrastive term still operates on the box branch exactly as
in the headline recipe. This is a structural change to \emph{where} the
domain-adversarial pressure is applied within the same SF-GRU
architecture, not a new architecture or a new dataset requirement.

Table~\ref{tab:context} shows this targeted variant matches the
untargeted headline recipe: both directions improve significantly over
the unadapted baseline, and neither direction differs from the
full-representation recipe by more than pooled standard deviation
(PIE~$\to$~JAAD: $0.617\pm0.020$ vs.\ $0.621\pm0.017$;
JAAD~$\to$~PIE: $0.637\pm0.055$ vs.\ $0.613\pm0.071$). If anything,
JAAD~$\to$~PIE's point estimate is slightly higher and its variance
is lower for the targeted variant, though this difference does not
clear pooled noise either. We read this as informative rather than as
a new headline number: a domain-adversarial term applied narrowly to
the diagnosed leakage source achieves the same benefit as applying it
broadly to the entire fused representation, touching a much smaller
part of the network to do so. This is consistent with our broader
reading (Section~\ref{sec:results-discussion}) that the recipe's value
lies in suppressing a specific, identifiable source of dataset-specific
signal rather than in a diffuse, whole-representation alignment effect
-- narrowing the intervention to that source loses nothing.

\begin{table}[t]
\centering
\caption{Context-targeted GRL (domain classifier reads only the
\texttt{local\_context} branch's hidden state) + box contrastive term,
$n=6$ seeds, corrected labels. Compare against Table~\ref{tab:fixes}'s
full-representation headline recipe.}
\label{tab:context}
\begin{tabular}{lcc}
\toprule
Method & PIE$\to$JAAD & JAAD$\to$PIE \\
\midrule
Baseline (Table~\ref{tab:baseline}) & $0.498\pm0.009$ & $0.402\pm0.138$ \\
GRL + contrastive (full repr., tuned) & $0.621\pm0.017$ & $0.613\pm0.071$ \\
GRL + contrastive (context-targeted) & $0.617\pm0.020$ & $0.637\pm0.055$ \\
\bottomrule
\end{tabular}
\end{table}

\subsection{Representational Separability Is Dissociated From Task Performance}
\label{sec:results-probe}

Table~\ref{tab:probe} reports linear and non-linear (MLP) probe
accuracy at distinguishing PIE from JAAD test-set representations
(Section~\ref{sec:method-probe}), for the unmodified baseline and for
several of the domain-adaptation configurations above. A trivial
control -- a probe given only each sample's track length, discarding
all learned features -- already achieves $84.6\%$ linear accuracy,
confirming PIE and JAAD differ in dataset-level statistics detectable
without any learned representation at all; all reported probe
accuracies should be read relative to this floor, not to chance
($50\%$).

The central finding is a dissociation between separability and task
performance. If domain-adversarial training worked by making the fused
representation literally domain-invariant, we would expect probe
accuracy to \emph{fall} as cross-dataset AUC \emph{rises}. We observe
the opposite: every configuration that improves cross-dataset AUC over
baseline (full GRL, GRL restricted to the context branch, GRL +
contrastive) also \emph{increases} probe accuracy over baseline, both
linear ($87.7\% \to 95.3$--$97.1\%$) and non-linear
($82.8\% \to 89.2$--$94.4\%$). Conversely, GRL + contrastive + mixup --
the one configuration in Table~\ref{tab:fixes} that \emph{reduces} AUC
relative to GRL + contrastive alone -- is also the one configuration
that \emph{reduces} probe accuracy toward baseline ($95.3\%\to88.3\%$
linear; $89.2\%\to83.3\%$ MLP). This pattern held across two
independent comparisons (the GRL family and the mixup control),
suggesting domain-adversarial training's benefit here is not explained
by the literal representational invariance it is nominally trained to
produce. We discuss a candidate mechanistic account -- that GRL
training acts as a generic regularizer on the fused representation,
incidentally raising probe-detectable structure while improving
task-relevant robustness through a different channel -- in
Section~\ref{sec:results-discussion}.

\begin{table}[t]
\centering
\caption{PIE-vs-JAAD probe accuracy on final fused hidden states (or,
where noted, a single branch's hidden state), averaged over 5 random
probe train/test splits. Higher accuracy = more separable, i.e.\ less
domain-invariant.}
\label{tab:probe}
\begin{tabular}{lcc}
\toprule
Condition & Linear acc. & MLP acc. \\
\midrule
Track-length only (trivial) & $0.846$ & $0.696\pm0.272$ \\
Baseline (fused) & $0.877\pm0.012$ & $0.828\pm0.022$ \\
GRL only (fused) & $0.962\pm0.017$ & $0.931\pm0.016$ \\
GRL, context branch only & $0.948\pm0.003$ & $0.894\pm0.038$ \\
Baseline, context branch only & $0.964\pm0.009$ & $0.851\pm0.026$ \\
GRL + contrastive (fused) & $0.953\pm0.008$ & $0.892\pm0.020$ \\
\ \ + mixup (fused) & $0.883\pm0.019$ & $0.833\pm0.013$ \\
\bottomrule
\end{tabular}
\end{table}

\subsection{Partial Replication on a Second Architecture}
\label{sec:results-crossattn}

The results above use a single fusion architecture (SF-GRU's stacked,
sequential concat-and-GRU fusion). To test whether the GRL + contrastive
recipe's benefit and the probe dissociation finding are specific to
this fusion mechanism or reflect something more general, we build the
same GRL + contrastive training recipe (Section~\ref{sec:method-dann})
on top of a structurally distinct architecture: box-anchored
cross-attention fusion (each modality independently GRU-encoded, the
box modality serving as an attention query over the concatenated other
modalities, a final GRU over the attended sequence), evaluated cross-
dataset for the first time here. At default contrastive
hyperparameters (no re-sweep for this architecture; see
Section~\ref{sec:method-dann}), Table~\ref{tab:crossattn} shows a
partial, direction-dependent replication. PIE~$\to$~JAAD improves over
the unmodified cross-attention baseline ($0.513\pm0.009 \to
0.559\pm0.027$, $n=6$, a difference that clears the two conditions'
pooled standard deviation), a smaller but real improvement in the same
direction as the stacked-GRU result. JAAD~$\to$~PIE does \emph{not}
replicate that improvement: the point estimate falls
($0.512\pm0.188 \to 0.446\pm0.063$), the opposite sign from the
stacked-GRU result on the same direction ($0.402\to0.613$), but the
baseline's own variance in this direction is large enough
($\text{std}=0.188$, matching the wide JAAD~$\to$~PIE variance noted
throughout this paper) that the drop does not clear pooled standard
deviation -- we can say GRL + contrastive fails to reproduce its
PIE~$\to$~JAAD-sized improvement here, but not confidently that it
makes this direction significantly worse. We report the point estimate
and its uncertainty together rather than rounding either away: the
fix's benefit is not confirmed to be architecture-independent, and a
practitioner should not assume it transfers to a new fusion design
without re-testing both directions at adequate $n$.

We probe both directions' cross-attention checkpoints
(Table~\ref{tab:probe-crossattn}), which reveals a more specific pattern
than "GRL + contrastive always increases separability." For
PIE~$\to$~JAAD, the direction that improves, separability rises
substantially over the unmodified baseline, both linearly
($0.844\pm0.006 \to 0.967\pm0.009$) and non-linearly
($0.846\pm0.000 \to 0.952\pm0.020$) -- matching every stacked-GRU
condition in Table~\ref{tab:probe}. For JAAD~$\to$~PIE, the direction
where GRL + contrastive fails to reproduce its improvement (its point
estimate falls, though not by more than the baseline's own noise;
Section~\ref{sec:results-crossattn}), separability barely moves at all
($0.845\pm0.004 \to 0.851\pm0.003$ linear; $0.828\pm0.033 \to
0.844\pm0.004$ non-linear) -- an order of magnitude smaller shift than
the PIE~$\to$~JAAD case. Separability and task performance are
therefore not simply decoupled in the sense of moving in opposite
directions; rather, the training procedure's measurable effect on the
representation tracks whether it helps the task at all. Where GRL +
contrastive meaningfully changes cross-dataset AUC (for better, on
stacked-GRU in both directions and cross-attention PIE~$\to$~JAAD), it
also meaningfully changes probe separability, always upward, never
downward; where it does not help (cross-attention JAAD~$\to$~PIE), it
leaves the representation close to where the unmodified baseline left
it, and the AUC change in that case is presumably driven by something
else in training (e.g.\ optimization dynamics specific to this
architecture and direction) rather than by any representational shift
this probe can detect. This refines, rather than weakens, our reading
in Section~\ref{sec:results-probe}: literal domain-invariance is not
what GRL + contrastive achieves in any condition we test, and the
magnitude of its (upward) effect on separability appears coupled to
whether it changes task behavior at all, not decoupled from it as a
simpler "always more separable" story would suggest.

\begin{table}[t]
\centering
\caption{Cross-attention probe accuracy, both transfer directions.
Baseline uses the checkpoint trained on the row's source dataset
(PIE-trained for the PIE$\to$JAAD row, JAAD-trained for JAAD$\to$PIE).}
\label{tab:probe-crossattn}
\begin{tabular}{lcc}
\toprule
Condition & Linear acc. & MLP acc. \\
\midrule
Baseline (PIE-trained) & $0.844\pm0.006$ & $0.846\pm0.000$ \\
GRL + contrastive, PIE$\to$JAAD & $0.967\pm0.009$ & $0.952\pm0.020$ \\
\midrule
Baseline (JAAD-trained) & $0.845\pm0.004$ & $0.828\pm0.033$ \\
GRL + contrastive, JAAD$\to$PIE & $0.851\pm0.003$ & $0.844\pm0.004$ \\
\bottomrule
\end{tabular}
\end{table}

\begin{table}[t]
\centering
\caption{Cross-attention architecture (box-anchored fusion), corrected
labels, $n=6$ seeds. Compare against Table~\ref{tab:fixes}'s
stacked-GRU numbers for the same comparison on SF-GRU.}
\label{tab:crossattn}
\begin{tabular}{lcc}
\toprule
Method & PIE$\to$JAAD & JAAD$\to$PIE \\
\midrule
Baseline (no adaptation) & $0.513\pm0.009$ & $0.512\pm0.188$ \\
GRL + contrastive & $0.559\pm0.027$ & $0.446\pm0.063$ \\
\bottomrule
\end{tabular}
\end{table}

\subsection{A Third Architecture: The Recipe Can Actively Hurt}
\label{sec:results-uncertainty}

We test a third, again structurally distinct architecture:
uncertainty-weighted fusion, in which each modality is independently
GRU-encoded and fused by a softmax over per-modality inverse-variance
weights (a learned confidence-weighted average) rather than sequential
concatenation (SF-GRU) or query-anchored attention (cross-attention).
This architecture is not an arbitrary third choice: it is the
strongest \emph{unadapted} baseline we find in an 8-way cross-dataset
architecture sweep (PIE~$\to$~JAAD AUC $0.585\pm0.009$, $n=3$,
exceeding SF-GRU's $0.498\pm0.009$ and cross-attention's
$0.513\pm0.009$ by a wide margin), making it the best-motivated
candidate for testing whether domain-adversarial training adds value
on top of an architecture that already generalizes comparatively well
without any adaptation. We apply the same GRL + contrastive recipe,
default hyperparameters, box-embedding contrastive tap on this
architecture's own independent box encoder output, $n=6$ seeds per
direction.

Table~\ref{tab:uncertainty} shows the recipe does not help here, and on
PIE~$\to$~JAAD actively \emph{hurts}: AUC falls from $0.585\pm0.009$ to
$0.499\pm0.072$, a drop that clears the two conditions' pooled standard
deviation and is the only case in this paper where GRL + contrastive
significantly \emph{under}performs an architecture's own unadapted
baseline. JAAD~$\to$~PIE also falls ($0.377\pm0.046 \to
0.351\pm0.060$), though this difference does not clear pooled standard
deviation. Combined with the cross-attention result
(Section~\ref{sec:results-crossattn}), only one of three architectures
we test (SF-GRU) shows the recipe helping in \emph{both} transfer
directions; cross-attention shows it helping in one and not
distinguishable from baseline in the other; uncertainty-weighted fusion
shows it significantly hurting in one and not distinguishable from
baseline (trending down) in the other. We revise our framing
accordingly: the earlier characterization of this recipe as producing
a ``partial, principled fix'' understates how architecture-contingent
it is. A more accurate summary is that GRL + contrastive is a
reproducible fix for SF-GRU's specific stacked-fusion topology, whose
mechanism (per Sections~\ref{sec:results-probe}
and~\ref{sec:results-crossattn}) does not appear to be literal
representational alignment, and which does not transfer as a
recommended recipe to the two alternative fusion mechanisms we
test -- one of which it actively harms on its strongest baseline.

\begin{table}[t]
\centering
\caption{Uncertainty-weighted fusion architecture, corrected labels.
Baseline is $n=3$ (from the 8-way architecture sweep); GRL + contrastive
is $n=6$. $\downarrow^{*}$ = drop clears pooled standard deviation;
$\downarrow$ = point estimate falls but not significantly at this $n$.}
\label{tab:uncertainty}
\begin{tabular}{lccc}
\toprule
Method & PIE$\to$JAAD & JAAD$\to$PIE & Verdict \\
\midrule
Baseline (no adaptation) & $0.585\pm0.009$ & $0.377\pm0.046$ & -- \\
GRL + contrastive & $0.499\pm0.072$ & $0.351\pm0.060$ & $\downarrow^{*}$/$\downarrow$ \\
\bottomrule
\end{tabular}
\end{table}

\subsection{A Pretrained Trajectory Representation Does Not Help}
\label{sec:results-pretrained-traj}

Gesnouin~et~al.~\cite{gesnouin2022} report that generic,
ImageNet/Sports1M-pretrained visual backbones generalize better under
PIE/JAAD domain shift than crossing-intention-specific visual
architectures -- motivating us to test the analogous question for the
\emph{motion} modality: does a box-trajectory representation pretrained
on a broader distribution than either dataset's own labeled training
pool generalize better cross-dataset than SF-GRU's from-scratch box
GRU? Lacking access to a large external trajectory corpus, we pretrain
a small Transformer encoder with a masked-frame reconstruction
objective (BERT-style: random frames of a box-delta sequence are
masked and predicted from context) on the pooled, \emph{unlabeled}
box-delta trajectories from both PIE's and JAAD's train and validation
splits combined -- a broader motion distribution than either dataset's
own labeled training pool alone, though a materially weaker substitute
for a genuine large-scale external corpus (e.g.\ Waymo Open Motion or
Argoverse~2), which we did not have local access to. We then freeze
this encoder and splice its per-timestep output into SF-GRU's box GRU
slot in place of the raw 4-dimensional box delta, training only the
downstream fusion (the box GRU, now consuming a 64-dimensional
embedding, and everything after it) under the same GRL + contrastive
recipe and default hyperparameters as Section~\ref{sec:results-dann}.

Table~\ref{tab:pretrained-traj} shows this does not help, and in fact
underperforms both the from-scratch GRL + contrastive result and, in
both directions, the plain unadapted baseline
(Table~\ref{tab:baseline}: $0.498$ and $0.402$). We do not read this as
evidence against a pretrained-backbone strategy in general, and we are
explicit about a confound in this experiment beyond corpus size alone:
our pretraining pool is drawn from PIE's and JAAD's \emph{own} train
and validation splits, the same two datasets used for the downstream
cross-dataset evaluation, whereas the visual-backbone analogy this
experiment is motivated by (Gesnouin~et~al.) concerns a backbone
pretrained on data with no relationship to either evaluation
domain. A self-supervised objective trained on a mixture of the source
and target domains has no pressure to discard whatever
dataset-identifying structure (frame rate, box-annotation convention,
scale) distinguishes the two -- indeed, our own probe methodology
(Section~\ref{sec:method-probe}) shows PIE and JAAD are separable from
nearly any representation of their inputs, including a single scalar
track length. It is therefore not clear whether this experiment tested
"does motion pretraining help cross-dataset transfer" or something
closer to "does an encoder pretrained to reconstruct a mixture of both
domains' own trajectories help," which is a different, and less
informative, question. Corpus size (a few hundred trajectories, vs.\
ImageNet's millions of images) compounds this: unlike a frozen VGG16
backbone pretrained on a large, source/target-independent corpus, a box
encoder pretrained on a small pool drawn from the very datasets under
study does not obviously carry information beyond what the
from-scratch, task-supervised box GRU already has access to, and the
added representation bottleneck (a frozen, non-task-tuned embedding
standing between the raw trajectory and the classification objective)
may cost more than the pretraining gives back at this scale. Isolating
which of these two factors -- domain composition or corpus size --
drives the negative result (for instance, by pretraining on PIE alone
and evaluating JAAD~$\to$~PIE, or vice versa, so the pretraining
corpus excludes the source domain entirely) is a direction for future
work we did not pursue here.

\begin{table}[t]
\centering
\caption{Frozen pretrained trajectory encoder spliced into the box
modality, GRL + contrastive recipe, default hyperparameters, corrected
labels, $n=6$ seeds. Compare against Table~\ref{tab:fixes}'s
from-scratch-box GRL + contrastive row (default hyperparameters:
$0.618\pm0.019$ / $0.583\pm0.073$).}
\label{tab:pretrained-traj}
\begin{tabular}{lcc}
\toprule
Method & PIE$\to$JAAD & JAAD$\to$PIE \\
\midrule
GRL + contrastive, pretrained traj.\ box & $0.489\pm0.019$ & $0.440\pm0.089$ \\
\bottomrule
\end{tabular}
\end{table}

\subsection{A Fine-Tuned Vision-Language Model Does Not Outperform SF-GRU}
\label{sec:results-vlm}

As a final, qualitatively different test of whether pretrained
world-knowledge helps this task, we LoRA-fine-tune Qwen2-VL-2B-Instruct
on a VQA-style (image-with-boxed-pedestrian, YES/NO) formulation of
crossing-intent, trained separately on each dataset and evaluated with
the same train-on-source/test-on-target protocol as every other result
in this paper. Unlike the from-scratch trajectory encoder of
Section~\ref{sec:results-pretrained-traj}, this backbone is pretrained
on a large, source/target-independent vision-language corpus, making it
a closer analogue to Gesnouin~et~al.'s generic-visual-backbone finding.
We report accuracy, precision, recall, and F1 rather than AUC: the
model is queried via greedy-decoded YES/NO generation, which yields a
hard prediction, not a calibrated probability, so these numbers are
\emph{not} directly comparable to the AUC figures used throughout the
rest of this paper.

In-domain accuracy is unremarkable relative to SF-GRU's in-domain AUC
(Table~\ref{tab:baseline}): $0.808$ (PIE, $n=224$) and, on JAAD, a
range of $0.647$--$0.824$ depending on which training checkpoint is
selected ($n=17$, see checkpoint-selection robustness below).
Cross-dataset, PIE~$\to$~JAAD accuracy is $0.581$ ($n=117$) with a
severe recall collapse (recall $0.154$: the model predicts
``not crossing'' for the large majority of samples regardless of
ground truth). For JAAD~$\to$~PIE, JAAD's validation split is only 17
sequences under \texttt{sample\_type=`beh'} -- the same structural
constraint noted throughout this paper for JAAD-as-source training
(Section~\ref{sec:results-discussion}) -- so selecting a single
``best'' fine-tuning checkpoint by validation accuracy on 17 samples is
itself high-variance (one flipped prediction moves accuracy by
$\approx\!6$ percentage points). Rather than report one
potentially-lucky checkpoint's cross-dataset number, we evaluate the
three best checkpoints by validation accuracy (training val.\ accuracy
$0.824$, $0.765$, $0.647$) independently on the full PIE test split
($n=637$ each). Table~\ref{tab:vlm-topk} shows this matters: the
checkpoint with the \emph{highest} training-time validation accuracy
is not the one with the best cross-dataset accuracy or F1, and the
spread across checkpoints (accuracy $0.281$--$0.474$) is large relative
to the differences we treat as meaningful elsewhere in this paper. We
report the mean and range across checkpoints rather than a single
point estimate for this reason.

Taken together, neither direction shows the fine-tuned VLM clearly
outperforming SF-GRU's from-scratch cross-dataset results (accuracy is
not directly comparable to our AUC numbers, but the severe recall
imbalance in several cells indicates the model is often not
discriminating well at all, independent of the choice of metric). We
read this as a second, qualitatively different negative result for the
``pretrained backbone rescues cross-dataset transfer'' hypothesis,
alongside Section~\ref{sec:results-pretrained-traj}'s trajectory-encoder
result -- though we note important differences in setup (a
single-image, VQA-style formulation with no explicit trajectory input,
vs.\ SF-GRU's multi-frame, multi-modality fusion) that mean this is not
a strictly controlled comparison, only a suggestive one.

\begin{table}[t]
\centering
\caption{VLM LoRA cross-dataset results, corrected labels. PIE$\to$JAAD
uses a single checkpoint (PIE's validation split, $n=224$, does not
have the small-sample selection problem JAAD's does). JAAD$\to$PIE
reports the three best JAAD-training checkpoints by validation accuracy
($n=17$ each), evaluated independently on PIE's full test split
($n=637$ each). Metrics are accuracy/F1/recall from greedy-decoded
YES/NO classification, not AUC.}
\label{tab:vlm-topk}
\begin{tabular}{lccc}
\toprule
Method & Acc. & F1 & Rec. \\
\midrule
PIE$\to$JAAD (single ckpt.) & $0.581$ & $0.246$ & $0.154$ \\
\midrule
JAAD$\to$PIE, rank 0 (best train val.) & $0.474$ & $0.394$ & $0.609$ \\
JAAD$\to$PIE, rank 1 & $0.422$ & $0.470$ & $0.911$ \\
JAAD$\to$PIE, rank 2 & $0.281$ & $0.439$ & $1.000$ \\
JAAD$\to$PIE, mean$\pm$std & $0.393\pm0.082$ & $0.434\pm0.031$ & $0.840\pm0.167$ \\
\bottomrule
\end{tabular}
\end{table}

\subsection{Split Sensitivity Is Not Explained by Video-Level Leakage}
\label{sec:results-leakage}

SF-GRU-family cross-dataset results are also sensitive to which
train/val/test partition is used within a dataset. Under JAAD's
official default split, our GRL + contrastive model reaches
PIE~$\to$~JAAD AUC $0.621\pm0.017$; under a random per-pedestrian
partition (\texttt{data\_split\_type=`random'}, ratios
$[0.5,0.4,0.1]$, freshly redrawn per run), the same method scores
$0.511\pm0.036$ -- a substantial drop. Because a per-pedestrian random
split does not guarantee that all pedestrians from a given video stay
in the same split, we hypothesized this drop was explained by
train/test video overlap: the default random split places $52\%$ of
JAAD pedestrians and (via a pedestrian-count coincidence) effectively
$100\%$ of PIE pedestrians in videos that appear on both sides of the
train/test boundary. We built a video-grouped variant of the random
split (Section~\ref{sec:method-probe}'s counterpart for splitting,
grouping all pedestrians in a video before partitioning, guaranteeing
zero video overlap) and re-ran the same method under it. Contrary to
our hypothesis, video-grouped-split AUC is \emph{lower}, not higher,
than the leaky pedestrian-level random split: $0.445\pm0.054$ ($n=3$)
vs.\ $0.511\pm0.036$. Video-level leakage is therefore not the
explanation for the default-split-vs-random-split gap; the true cause
remains open. We report this as a negative result rather than
discarding it, since it rules out what seemed the most likely
explanation and narrows the space of remaining hypotheses (e.g.\ the
default splits may differ systematically in scene composition or
pedestrian-behavior distribution in a way a random partition does not
control for).

\subsection{Summary of Everything Attempted}
\label{sec:results-summary}

Table~\ref{tab:summary} collects every fix attempt in this paper in one
place, since the individual results are otherwise spread across eight
subsections written around different hypotheses, with a verdict
computed against pooled standard deviation rather than by eye. Of the
thirteen non-baseline interventions, only the domain-adversarial family
on SF-GRU (GRL alone, GRL + contrastive full-representation, GRL +
contrastive context-targeted) improves PIE~$\to$~JAAD \emph{and}
JAAD~$\to$~PIE by more than pooled noise in both directions. The
geometry- and normalization-based fixes (rows 2--5) all significantly
\emph{hurt} PIE~$\to$~JAAD and are statistically indistinguishable
from baseline on JAAD~$\to$~PIE -- not the symmetric two-direction harm
a first read of the raw numbers might suggest, since JAAD~$\to$~PIE's
wide baseline variance (row 1; Section~\ref{sec:results-discussion})
absorbs differences that would be significant in the other direction.
The context-targeted variant (row 9) is statistically indistinguishable
from the full-representation recipe (row 8) in both directions
(Section~\ref{sec:results-context}) -- narrowing the domain-adversarial
pressure to the diagnosed leakage source loses nothing. Synthetic speed
and the pretrained trajectory representation (rows 11--12) both
significantly hurt \emph{both} directions relative to the tuned
headline recipe; mixup (row 10) is statistically indistinguishable from
it in both directions despite its lower point estimate. Applied to the
two alternative fusion architectures (rows 13--14), the recipe's record
is mixed-to-negative: on cross-attention it significantly helps
PIE~$\to$~JAAD and is statistically indistinguishable from its own
baseline (not significantly worse) on JAAD~$\to$~PIE; on
uncertainty-weighted fusion -- the strongest unadapted baseline of any
architecture we test -- it significantly \emph{hurts} PIE~$\to$~JAAD
and trends down (not significant) on JAAD~$\to$~PIE. Only SF-GRU shows
the recipe helping in both directions, whether applied broadly or
narrowly. Table~\ref{tab:summary} reports
AUC throughout and so excludes the VLM LoRA results
(Section~\ref{sec:results-vlm}, which use greedy-decoded accuracy/F1
rather than a probability score); those results are a further,
qualitatively different negative data point -- a large pretrained
vision-language backbone also does not clearly outperform SF-GRU's
from-scratch cross-dataset numbers -- that this table's AUC-only format
cannot represent alongside the rest. We read this table as the
paper's most concise evidence for its central claim: cross-dataset
collapse in this model family is not fixed by any single targeted
intervention we tried except domain-adversarial training on SF-GRU
specifically, and that fix is not just partial but architecture-
contingent, actively harmful on at least one alternative architecture's
strongest baseline.

\begin{table}[t]
\centering
\caption{Every intervention attempted in this paper, both transfer
directions, corrected labels. SF-GRU (stacked fusion) unless noted; row
4 (homography) is the metric ground-plane reprojection of
Section~\ref{sec:method-geometry}, row 5 drops the box modality
entirely. Verdict is $\uparrow$/$\downarrow$/$\approx$ relative to the
stated baseline, where $\approx$ means the difference is smaller than
the two conditions' pooled standard deviation (i.e.\ not
distinguishable at this $n$) -- rows 2--9 compare against the unadapted
baseline (row 1); rows 10--12 compare against GRL + contrastive tuned
(row 8); rows 13--14 compare against their own architecture's unadapted
baseline (Tables~\ref{tab:crossattn} and~\ref{tab:uncertainty}).}
\label{tab:summary}
\small
\begin{tabular}{lccc}
\toprule
Intervention & PIE$\to$JAAD & JAAD$\to$PIE & Verdict \\
\midrule
1. Baseline (no fix) & $0.498\pm0.009$ & $0.402\pm0.138$ & -- \\
2. Scale-norm & $0.437\pm0.025$ & $0.536\pm0.063$ & $\downarrow$/$\approx$ \\
3. Height-norm & $0.390\pm0.017$ & $0.554\pm0.104$ & $\downarrow$/$\approx$ \\
4. Homography & $0.406\pm0.011$ & $0.544\pm0.053$ & $\downarrow$/$\approx$ \\
5. No-box & $0.385\pm0.009$ & $0.432\pm0.085$ & $\downarrow$/$\approx$ \\
6. Contrastive only & $0.558\pm0.028$ & $0.450\pm0.107$ & $\uparrow$/$\approx$ \\
7. GRL only & $0.608\pm0.030$ & $0.587\pm0.069$ & $\uparrow$/$\uparrow$ \\
8. GRL+contrastive (tuned) & $\mathbf{0.621\pm0.017}$ & $\mathbf{0.613\pm0.071}$ & $\uparrow$/$\uparrow$ \\
9. GRL+contrastive (context-targeted) & $0.617\pm0.020$ & $0.637\pm0.055$ & $\uparrow$/$\uparrow$ \\
10.\ + mixup & $0.580\pm0.077$ & $0.537\pm0.079$ & $\approx$/$\approx$ \\
11.\ + synth.\ speed & $0.602\pm0.001$ & $0.430\pm0.072$ & $\downarrow$/$\downarrow$ \\
12.\ + pretrained traj.\ & $0.489\pm0.019$ & $0.440\pm0.089$ & $\downarrow$/$\downarrow$ \\
13.\ Cross-attn.\ & $0.559\pm0.027$ & $0.446\pm0.063$ & $\uparrow$/$\approx$ \\
14.\ Uncertainty fusion & $0.499\pm0.072$ & $0.351\pm0.060$ & $\downarrow$/$\approx$ \\
\bottomrule
\end{tabular}
\end{table}

\subsection{Discussion}
\label{sec:results-discussion}

Four results together shape our interpretation. First, geometry- and
normalization-based fixes are uniformly ineffective
(Section~\ref{sec:results-geometry}), ruling out the most
straightforward explanation (differing camera setups) for
cross-dataset collapse. Second, the modality-ablation result
(Section~\ref{sec:results-ablation}) and the probe results
(Section~\ref{sec:results-probe}) together point away from a
single-modality, single-cause account: visual context matters most to
zero out, yet the mechanism that helps most (GRL) does not work by
suppressing a detectable dataset signal -- it \emph{increases} that
signal's detectability while improving task performance. We read this
as evidence that GRL's practical benefit in this setting is closer to
a generic regularization effect on the fused representation (e.g.\
reducing reliance on dataset-specific spurious correlations the main
task loss alone would otherwise fit) than to the literal
representation-alignment mechanism it is designed around, echoing
similar decoupling results reported for DANN-family methods outside
this domain. Third, the cross-attention replication
(Section~\ref{sec:results-crossattn}) sharpens this reading further:
probing both transfer directions shows the separability increase is
not a fixed side effect of the training procedure, but tracks whether
GRL + contrastive changes task behavior at all -- large where it helps
(PIE~$\to$~JAAD on both architectures, JAAD~$\to$~PIE on stacked-GRU)
and negligible where it does not (JAAD~$\to$~PIE on cross-attention). A
representation-alignment account predicts a consistent effect on
separability regardless of whether AUC improves; what we observe -- the
separability shift and the AUC shift rising and falling together in
magnitude, even though the AUC shift is not always positive -- is more
consistent with GRL + contrastive perturbing the representation only
when it meaningfully engages with training at all, and that engagement
sometimes helping and sometimes hurting the task depending on
architecture-specific training dynamics. The uncertainty-weighted
fusion result (Section~\ref{sec:results-uncertainty}) is the sharpest
data point for this view: applied to the architecture with the
\emph{best} unadapted baseline of any we test, the same recipe
significantly \emph{degrades} it. If GRL + contrastive's benefit came
from a mechanism intrinsic to domain-adversarial training working
correctly, we would not expect it to reliably help one architecture
and reliably hurt another; a generic-regularizer account, where the
training procedure's effect on a given architecture's specific
optimization landscape can go either way, accommodates this more
naturally than an account in which the method's success is attributed
to achieving its nominal objective. Fourth, the persistent, large
JAAD~$\to$~PIE variance across
nearly every method in Tables~\ref{tab:fixes},~\ref{tab:crossattn},
and~\ref{tab:uncertainty}
deserves an explicit alternative explanation before we attribute it to
the benchmark's scene composition: JAAD's default-split training set
has 103 labeled sequences, against PIE's 796 -- an 8-fold difference --
so JAAD-as-source runs are trained on a much smaller, and therefore
individually noisier, sample than PIE-as-source runs, quite apart from
any property of \emph{which} scenes each dataset contains. A smaller
source training set producing a noisier decision boundary, and hence
higher-variance transfer to any target, is a mundane statistical
explanation that competes directly with a benchmark-composition
account, and our experiments do not distinguish between them: every
JAAD-as-source condition in this paper confounds ``JAAD'' with
``the smaller-training-set direction.'' The counter-hypothesis
split-sensitivity result (Section~\ref{sec:results-leakage}) is
consistent with either account. We flag this rather than resolve it: a
clean test would hold training-set size fixed (e.g.\ subsampling PIE's
source training set to JAAD's size) and check whether PIE~$\to$~JAAD
variance rises to match, which we did not run. Together, this and the
split-sensitivity result suggest that some of
this task's difficulty is a property of the benchmark itself --
scene composition, training-set size, or both -- rather than purely a
property of any model or training method evaluated on it. We view
robust cross-dataset pedestrian intent prediction as still substantially
unsolved, with our contribution best read as (a) a fix that is
reproducible and genuinely useful for one specific architecture
(SF-GRU's stacked fusion) but is not a general-purpose recipe --
tested on two further, structurally distinct architectures, it helps
partially on one and significantly hurts the other's best baseline in
at least one direction -- together with a mechanistic account
(Sections~\ref{sec:results-probe}--\ref{sec:results-uncertainty}) of
why its effect is architecture-contingent rather than a property of
whether it achieves genuine domain-invariance, and (b) a set of
ruled-out hypotheses -- geometry, video-level leakage, a missing speed
modality, and a locally-pretrained trajectory representation -- that we
hope narrows the search space for future work more than our single
positive result advances it.


% conclusion.tex

\section{Conclusion}
\label{sec:conclusion}

We set out to understand why SF-GRU-family pedestrian crossing-intention
models fail to generalize across PIE and JAAD, and found more than one
answer was needed. Part of the difficulty was a silent, orthogonal data
bug in JAAD's default configuration, unrelated to cross-dataset transfer
itself but affecting any result built on the unaudited default; we
identify, fix, and re-confirm our findings hold under the correction.
The remaining, genuine cross-dataset collapse is not explained by camera
geometry, and the closely related phenomenon of split-sensitivity is
not explained by video-level train/test leakage -- both plausible,
natural hypotheses that we falsify directly rather than merely fail to
support. A domain-adversarial and contrastive training recipe, applied
to the fused representation and motivated by a modality-zeroing
ablation, produces a real, seed-confirmed improvement in both transfer
directions on SF-GRU, but leaves cross-dataset performance well below
in-domain levels; we consider this residual gap, not the improvement,
to be the more important number to report honestly. A
representation-separability probe, deliberately independent of the
domain classifier used during training, shows this recipe's benefit is
dissociated from the literal domain-invariance it is designed to
induce -- the representations that help most are, if anything,
\emph{more} separable by an external probe, not less -- suggesting its
practical effect is closer to a generic regularizer on the fused
representation than to representation alignment per se. Testing the
same recipe on two further, structurally distinct fusion architectures
supports this reading further and tempers any claim of a general-
purpose fix: on cross-attention fusion the improvement replicates only
partially, and on uncertainty-weighted fusion -- the strongest
unadapted baseline architecture we find -- the same recipe
significantly \emph{degrades} performance. We take from this that
domain-adversarial training's value here is real but architecture-
specific, not evidence that it achieves anything close to its nominal
objective of representational invariance regardless of architecture.

We draw two broader conclusions. First, cross-dataset evaluation should
be treated as standard practice for pedestrian crossing-intention
architectures, not an optional follow-up: every architecture we survey
in Section~\ref{sec:related} published since Gesnouin~et~al.'s original
finding~\cite{gesnouin2022} reports only in-domain results, and our own
in-domain numbers (Table~\ref{tab:baseline}) look entirely unremarkable
next to a near-chance cross-dataset collapse the same checkpoint
exhibits. In-domain performance on PIE or JAAD alone is not informative
about whether a proposed architecture has learned anything that
transfers. Second, we think auditing benchmark and pipeline details
that are easy to get silently wrong -- data-loader defaults, split
construction, the provenance of a supposedly domain-invariant
representation -- is undersupplied relative to the rate at which new
architectures are proposed for these benchmarks, and that the negative
results in this paper (ruled-out geometry, ruled-out leakage, a
dissociated invariance mechanism) are worth reporting even though none
of them, individually, is the paper's headline number.

\textbf{Limitations.} Our experiments center on SF-GRU's stacked fusion;
we additionally test our headline recipe on two further, structurally
distinct architectures (box-anchored cross-attention,
Section~\ref{sec:results-crossattn}, and uncertainty-weighted fusion,
Section~\ref{sec:results-uncertainty}). The probe dissociation --
higher separability wherever the recipe changes task behavior at all --
replicates on both, but the \emph{fix itself} does not: cross-attention
shows a partial, one-direction improvement, and uncertainty-weighted
fusion shows the recipe significantly harming its (strongest, of any
architecture we test) unadapted baseline. We therefore have reasonable
evidence the dissociation mechanism is not an SF-GRU-specific artifact,
but active evidence \emph{against} the fix being architecture-general,
not merely an absence of evidence for it. We did not test transformer-
or Mamba-based architectures at all, so whether our findings extend to
substantially larger or differently-inductive-biased fusion mechanisms
remains open, and a practitioner adopting this recipe on a new
architecture should re-validate both directions rather than assume
transfer. We evaluate on two datasets; both are
pedestrian-crossing-specific,
relatively small by modern standards, and share enough protocol
lineage (both from the same research group) that a third, more
independently sourced dataset (e.g.\ TITAN) would strengthen any claim
that our findings reflect a general property of this task rather than
a property specific to the PIE/JAAD pair. Our camera-geometry
falsification relies on a researched, not measured, camera calibration
for JAAD, since JAAD does not publish one; we mitigate this with a
sensitivity sweep over plausible parameter ranges but cannot rule out
that a parameter setting outside our swept range behaves differently.
Finally, our synthesized JAAD speed modality is a coarse optical-flow
proxy, not a calibrated measurement, and its negative result should be
read as evidence about this particular proxy rather than as a general
claim that a speed signal cannot help JAAD cross-dataset transfer.

We release our corrected-labeling configuration, falsification
experiments, and probe methodology in the hope that they are directly
reusable by future work evaluating new architectures on PIE and JAAD,
independent of whether that work adopts our specific training recipe.


\bibliographystyle{IEEEtran}
\bibliography{refs}

\end{document}
